% ============================================================
% Annexe : documentation open source de RevMine
% ============================================================

\section{Guide d'installation}

RevMine peut etre execute localement avec Docker Compose. Les prerequis minimaux sont Git,
Docker Desktop ou Docker Engine avec le plugin Compose, et un espace memoire suffisant pour
executer les bases PostgreSQL, Kafka, MinIO, Grafana/Loki et les services applicatifs. Pour un
developpement hors conteneur, Node.js est necessaire pour le frontend, Python 3.11 pour les
services Python et Go pour le service de notification.

La configuration de depart est fournie par \texttt{.env.example}. Le fichier doit etre copie
vers \texttt{.env}, puis complete avec les secrets de l'environnement local.

\begin{lstlisting}[caption={Initialisation de l'environnement local RevMine},label={lst:annexe-install-env},breaklines=true]
cp .env.example .env
# renseigner SECRET_KEY, ENCRYPTION_KEY, POSTGRES_USER,
# POSTGRES_PASSWORD, les credentials OAuth et les tokens LLM si besoin
docker compose up --build
\end{lstlisting}

Les variables les plus sensibles sont \texttt{SECRET\_KEY}, \texttt{ENCRYPTION\_KEY}, les
identifiants PostgreSQL, les credentials OAuth GitHub/GitLab/Google, les parametres MinIO
(\texttt{MINIO\_ROOT\_USER}, \texttt{MINIO\_ROOT\_PASSWORD}, \texttt{MINIO\_BUCKET\_NAME}) et les
cles eventuelles pour OpenRouter. La cle \texttt{ENCRYPTION\_KEY} protege les tokens Git stockes
par le service de configuration; elle doit donc etre traitee comme un secret d'application.

\begin{table}[H]
\centering
\small
\caption{Acces locaux principaux}
\label{tab:annexe-acces-locaux}
\begin{tabular}{|>{\raggedright\arraybackslash}p{4cm}
                |>{\raggedright\arraybackslash}p{5cm}
                |>{\raggedright\arraybackslash}p{5cm}|}
\hline
\textbf{Composant} & \textbf{URL locale} & \textbf{Usage} \\
\hline
Frontend & \texttt{http://localhost:5173} & Interface React/Vite. \\
\hline
API Gateway & \texttt{http://localhost:8000} & Point d'entree REST principal. \\
\hline
MinIO Console & \texttt{http://localhost:9001} & Consultation des objets JSON et CSV. \\
\hline
Grafana & \texttt{http://localhost:3001} & Dashboards et requetes Loki. \\
\hline
Loki & \texttt{http://localhost:3100} & API de stockage des logs. \\
\hline
SonarQube & \texttt{http://localhost:9090} & Pile separee via \texttt{docker-compose.sonar.yml}. \\
\hline
\end{tabular}
\end{table}

Si l'inference locale est utilisee, un modele Ollama doit etre disponible dans le conteneur
\texttt{ollama-service}. Le service LLM peut aussi utiliser OpenRouter lorsque la cle
correspondante est fournie dans l'environnement.

\section{Guide developpeur}

Le depot est organise autour d'une separation claire entre frontend, gateway et microservices :

\begin{itemize}
    \item \texttt{frontend/} : application React/Vite, routes protegees, services Axios,
          visualisations ECharts et contexte de notifications.
    \item \texttt{backend/api-gateway/} : gateway FastAPI, routage, CORS, introspection JWT,
          propagation de \texttt{X-User-ID} et rate limiting.
    \item \texttt{backend/services/authentication/} : comptes utilisateurs, JWT SimpleJWT,
          refresh token, introspection et OAuth GitHub/GitLab/Google.
    \item \texttt{backend/services/configuration/} : workspaces, tokens chiffres, import de depots
          GitHub/GitLab et reponse aux demandes Kafka de token.
    \item \texttt{backend/services/collection/} : plans de collecte, collecte GitHub/GitLab,
          stockage MinIO, nettoyage, CSV et evenements Kafka.
    \item \texttt{backend/services/analyze/} : datasets, catalogue de metriques, calculs
          analytiques, taches Celery et modules DevOps Kanban/CI-CD.
    \item \texttt{backend/services/llm/} : normalisation JSON des demandes utilisateur via
          Ollama ou OpenRouter.
    \item \texttt{backend/services/notification/} : service Go/Fiber consommant Kafka et diffusant
          les notifications REST/WebSocket.
\end{itemize}

Les tests peuvent etre lances service par service. Le pipeline GitLab utilise ces memes
separations afin de ne tester que les services modifies.

\begin{lstlisting}[caption={Commandes de test par famille de services},label={lst:annexe-tests},breaklines=true]
# Services Python
cd backend/api-gateway && python -m pytest
cd backend/services/authentication && python -m pytest
cd backend/services/configuration && python -m pytest
cd backend/services/collection && python -m pytest
cd backend/services/analyze && python -m pytest
cd backend/services/llm && python -m pytest

# Service Go
cd backend/services/notification && go test ./...

# Frontend
cd frontend && npm test
cd frontend && npm run cypress:run
\end{lstlisting}

Pour ajouter un endpoint dans un service Django/DRF, la demarche recommandee est d'ajouter le
serializer ou le cas d'usage, d'exposer la vue DRF, de declarer la route dans \texttt{urls.py},
puis de couvrir le comportement par des tests \texttt{pytest-django}. Pour un endpoint FastAPI,
le changement doit etre accompagne d'un schema de requete/reponse Pydantic et d'un test HTTP.

Pour ajouter une metrique, il faut completer le catalogue \texttt{MetricDefinition}, implementer
la fonction de calcul dans le moteur de metriques, declarer les contraintes de colonnes et d'axes
compatibles, puis ajouter le mapping frontend lorsque la visualisation necessite un traitement
specifique.

\section{Guide API}

La documentation API est exposee differemment selon les technologies utilisees. Les URLs ci-dessous
sont des acces de developpement directs aux services; dans l'architecture cible, les appels
applicatifs passent par l'API Gateway.

\begin{table}[H]
\centering
\small
\caption{Documentation API disponible par service}
\label{tab:annexe-doc-api}
\begin{tabular}{|>{\raggedright\arraybackslash}p{4cm}
                |>{\raggedright\arraybackslash}p{5cm}
                |>{\raggedright\arraybackslash}p{5cm}|}
\hline
\textbf{Service} & \textbf{Documentation} & \textbf{Observation} \\
\hline
API Gateway & \texttt{/docs}, \texttt{/redoc}, \texttt{/openapi.json} & Documentation FastAPI du gateway. \\
\hline
Authentication & \texttt{/api/docs/}, \texttt{/api/redoc/}, \texttt{/api/schema/} & Documentation DRF Spectacular. \\
\hline
Configuration & \texttt{/api/docs/}, \texttt{/api/redoc/}, \texttt{/api/schema/} & Workspaces, tokens et depots. \\
\hline
Collection & \texttt{/api/docs/}, \texttt{/api/redoc/}, \texttt{/api/schema/} & Plans, nettoyage et datasets collectes. \\
\hline
LLM & \texttt{/docs}, \texttt{/redoc}, \texttt{/openapi.json} & Documentation FastAPI du parser LLM. \\
\hline
Analyze & Non declare dans \texttt{analyze\_service/urls.py} & Routes disponibles sous \texttt{/api/analysis/}. \\
\hline
Notification & Non declare & API Go/Fiber REST et WebSocket documentee par le code. \\
\hline
\end{tabular}
\end{table}

\section{Guide des metriques}

Le catalogue des metriques est porte par le service d'analyse, principalement autour du modele
\texttt{MetricDefinition}, des migrations de donnees du module \texttt{analytics} et du moteur de
calcul situe dans le domaine analytique. Ce catalogue separe le nom fonctionnel de la metrique,
sa categorie, ses contraintes de colonnes et la fonction de calcul appelee au moment de generer
une analyse.

\begin{table}[H]
\centering
\small
\caption{Exemples de metriques presentes dans le catalogue RevMine}
\label{tab:annexe-metrics}
\begin{tabular}{|>{\raggedright\arraybackslash}p{3.8cm}
                |>{\raggedright\arraybackslash}p{4.2cm}
                |>{\raggedright\arraybackslash}p{6cm}|}
\hline
\textbf{Metrique} & \textbf{Famille} & \textbf{Objectif analytique} \\
\hline
\texttt{lead\_time\_distribution} & Revue de code & Mesurer le temps entre creation et integration d'une PR/MR. \\
\hline
\texttt{code\_churn} & Evolution du code & Suivre les lignes ajoutees, supprimees et modifiees. \\
\hline
\texttt{entropy\_analysis} & Dispersion historique & Evaluer la concentration des changements sur fichiers ou modules. \\
\hline
\texttt{top\_reviewers} & Engagement reviewers & Identifier les reviewers les plus impliques. \\
\hline
\texttt{time\_to\_first\_review} & Revue de code & Quantifier le delai avant le premier retour. \\
\hline
\texttt{review\_duration}, \texttt{cycle\_time} & Flux de revue & Mesurer la duree de revue et le cycle complet. \\
\hline
\texttt{kanban\_lead\_time}, \texttt{kanban\_cycle\_time}, \texttt{kanban\_throughput} & Kanban & Analyser le flux des issues et la capacite de livraison. \\
\hline
\texttt{cicd\_success\_rate}, \texttt{cicd\_build\_duration}, \texttt{cicd\_mttr} & CI/CD & Suivre la fiabilite et la performance des pipelines. \\
\hline
\end{tabular}
\end{table}

\section{Guide d'extension}

RevMine est concu pour etre etendu par ajout progressif de collecteurs, de metriques ou de vues
frontend. L'architecture microservices limite l'impact d'une extension : un nouveau collecteur
peut etre integre dans le service de collecte ou dans le module DevOps du service d'analyse,
tandis qu'une nouvelle metrique peut etre ajoutee au catalogue sans modifier l'authentification
ni la configuration des workspaces.

Deux extensions sont particulierement naturelles :

\begin{itemize}
    \item \textbf{Analyse avancee des tableaux Kanban} : enrichir la collecte des issues,
          transitions, labels et colonnes pour produire davantage d'indicateurs de flux.
    \item \textbf{Analyse des builds et pipelines CI} : exploiter les jobs, statuts,
          durees, echecs et reprises pour relier qualite de revue, stabilite CI et performance de
          livraison.
\end{itemize}

Cette extensibilite repose sur quatre points : l'isolation des services, le stockage brut dans
MinIO, le catalogue de metriques du service d'analyse et la separation frontend entre parcours,
services API et composants de visualisation.
